%%%%%%%%%%%%%%%%%%%%%%%%%%%%%%%%%%%%%%%%%
% Wenneker Assignment
% LaTeX Template
% Version 2.0 (12/1/2019)
%
% This template originates from:
% http://www.LaTeXTemplates.com
%
% Authors:
% Vel (vel@LaTeXTemplates.com)
% Frits Wenneker
%
% License:
% CC BY-NC-SA 3.0 (http://creativecommons.org/licenses/by-nc-sa/3.0/)
% 
%%%%%%%%%%%%%%%%%%%%%%%%%%%%%%%%%%%%%%%%%

%----------------------------------------------------------------------------------------
%	PACKAGES AND OTHER DOCUMENT CONFIGURATIONS
%----------------------------------------------------------------------------------------

\documentclass[11pt]{scrartcl} % Font size

%%%%%%%%%%%%%%%%%%%%%%%%%%%%%%%%%%%%%%%%%
% Wenneker Assignment
% Structure Specification File
% Version 2.0 (12/1/2019)
%
% This template originates from:
% http://www.LaTeXTemplates.com
%
% Authors:
% Vel (vel@LaTeXTemplates.com)
% Frits Wenneker
%
% License:
% CC BY-NC-SA 3.0 (http://creativecommons.org/licenses/by-nc-sa/3.0/)
% 
%%%%%%%%%%%%%%%%%%%%%%%%%%%%%%%%%%%%%%%%%

%----------------------------------------------------------------------------------------
%	PACKAGES AND OTHER DOCUMENT CONFIGURATIONS
%----------------------------------------------------------------------------------------

\usepackage{amsmath, amsfonts, amsthm} % Math packages

\usepackage{listings} % Code listings, with syntax highlighting
\usepackage[outputdir=out]{minted}

\usepackage[english]{babel} % English language hyphenation

\usepackage{graphicx} % Required for inserting images
\graphicspath{{Figures/}{./}} % Specifies where to look for included images (trailing slash required)

\usepackage{booktabs} % Required for better horizontal rules in tables

\numberwithin{equation}{section} % Number equations within sections (i.e. 1.1, 1.2, 2.1, 2.2 instead of 1, 2, 3, 4)
\numberwithin{figure}{section} % Number figures within sections (i.e. 1.1, 1.2, 2.1, 2.2 instead of 1, 2, 3, 4)
\numberwithin{table}{section} % Number tables within sections (i.e. 1.1, 1.2, 2.1, 2.2 instead of 1, 2, 3, 4)

\setlength\parindent{0pt} % Removes all indentation from paragraphs

\usepackage{enumitem} % Required for list customisation
\setlist{noitemsep} % No spacing between list items

%----------------------------------------------------------------------------------------
%	DOCUMENT MARGINS
%----------------------------------------------------------------------------------------

\usepackage{geometry} % Required for adjusting page dimensions and margins

\geometry{
	paper=a4paper, % Paper size, change to letterpaper for US letter size
	top=2.5cm, % Top margin
	bottom=3cm, % Bottom margin
	left=3cm, % Left margin
	right=3cm, % Right margin
	headheight=0.75cm, % Header height
	footskip=1.5cm, % Space from the bottom margin to the baseline of the footer
	headsep=0.75cm, % Space from the top margin to the baseline of the header
	%showframe, % Uncomment to show how the type block is set on the page
}

%----------------------------------------------------------------------------------------
%	FONTS
%----------------------------------------------------------------------------------------

\usepackage[utf8]{inputenc} % Required for inputting international characters
\usepackage[T1]{fontenc} % Use 8-bit encoding

\usepackage{fourier} % Use the Adobe Utopia font for the document

%----------------------------------------------------------------------------------------
%	SECTION TITLES
%----------------------------------------------------------------------------------------

\usepackage{sectsty} % Allows customising section commands

\sectionfont{\vspace{6pt}\centering\normalfont\scshape} % \section{} styling
\subsectionfont{\normalfont\bfseries} % \subsection{} styling
\subsubsectionfont{\normalfont\itshape} % \subsubsection{} styling
\paragraphfont{\normalfont\scshape} % \paragraph{} styling

%----------------------------------------------------------------------------------------
%	HEADERS AND FOOTERS
%----------------------------------------------------------------------------------------

\usepackage{scrlayer-scrpage} % Required for customising headers and footers

\ohead*{} % Right header
\ihead*{} % Left header
\chead*{} % Centre header

\ofoot*{} % Right footer
\ifoot*{} % Left footer
\cfoot*{\pagemark} % Centre footer
 % Include the file specifying the document structure and custom commands

%----------------------------------------------------------------------------------------
%	TITLE SECTION
%----------------------------------------------------------------------------------------

\title{	
	\normalfont\normalsize
	\begin{center}
		\begin{minipage}[c]{0.2\textwidth}
			\textsc{\Large SZ-Ybbs}
		\end{minipage}%
		\begin{minipage}[c]{0.1\textwidth}
			\includegraphics[width=\textwidth]{LogoITHTL_white.pdf}
		\end{minipage}
	\end{center}
	\vspace{10pt} % Whitespace
	\rule{\linewidth}{0.5pt}\\ % Thin top horizontal rule
	\vspace{20pt} % Whitespace
	{\huge TLS}\\ % The assignment title
	\vspace{12pt} % Whitespace
	\rule{\linewidth}{2pt}\\ % Thick bottom horizontal rule
	\vspace{12pt} % Whitespace
}

\author{\LARGE Erber Jakob, Freunberger Raphael} % Your name

\date{\normalsize\today} % Today's date (\today) or a custom date

\begin{document}

\maketitle % Print the title

\section{TLS 1.2/1.3 Overview}

Transport Layer Security (TLS) is a cryptographic protocol used for providing secure communications over an insecure network, like the internet. TSL provides security, confidentiality authenticate and integrity. TSL builds on top of the now-deprecated Secure Sockets Layer (SSL) protocol and does not clearly fit one OSI-layer. Depending on who you ask TLS can be a layer 4, 6 or 7 protocol. The protocol uses TCP, but there is also a UDP equivalent called Datagram Transport Layer Security (DTLS). DTLS is sometimes used for VPNs but in general used less than TSL. TSL has a wide range of use cases, from email to voice over IP or web browsing.
\\\\
To keep its security promise TLS has to keep being developed and new versions get released as the older versions become insecure and get deprecated.

\begin{table}[h]
    \centering
    \begin{tabular}{|l|c|l|}
        \hline
        \textbf{Version} & \textbf{Year} & \textbf{Status} \\ \hline
        TLS 1.0 (insecure) & 1999 & Deprecated in 2021 (RFC 8996) \\ \hline
        TLS 1.1 (insecure) & 2006 & Deprecated in 2021 (RFC 8996) \\ \hline
        TLS 1.2 & 2008 & In use since 2008 \\ \hline
        TLS 1.3 & 2018 & In use since 2018 \\ \hline
    \end{tabular}
    \caption{TLS Versions and Status}
    \label{tab:tls_versions}
\end{table}

\subsection{TLS supported ciphers}
\label{sec:tls_supported_ciphers}

These are the supported ciphers for the different TLS versions. The table shows which ciphers are secure and which are not. The information is taken from wikipedia. This table will be used later to compare the security of the different TLS versions and supported ciphers.

\begin{table}[h!]
    \centering
    \renewcommand{\arraystretch}{1.3}
    \begin{tabular}{|p{2cm}|p{3cm}|p{2cm}|p{2cm}|p{2cm}|p{2cm}|}
        \hline
        \textbf{Type} & \textbf{Algorithm} & \textbf{TLS 1.0} & \textbf{TLS 1.1} & \textbf{TLS 1.2} & \textbf{TLS 1.3} \\ \hline
        
        \multirow{6}{*}{Block} & AES GCM & — & — & Secure & Secure \\ \cline{2-6}
         & AES CCM & — & — & Secure & Secure \\ \cline{2-6}
         & AES CBC & Insecure & Depends on mitigations & Depends on mitigations & — \\ \cline{2-6}
         & Camellia GCM & — & — & Secure & — \\ \cline{2-6}
         & Camellia CBC & Insecure & Depends on mitigations & Depends on mitigations & — \\ \cline{2-6}
         & ARIA GCM & — & — & Secure & — \\ \cline{2-6}
         & ARIA CBC & — & Depends on mitigations & Depends on mitigations & — \\ \cline{2-6}
         & SEED CBC & Insecure & Depends on mitigations & Depends on mitigations & — \\ \cline{2-6}
         & 3DES EDE CBC & Insecure & Insecure & Insecure & — \\ \cline{2-6}
         & GOST R 34.12-2015 Magma CTR & — & — & Insecure & — \\ \cline{2-6}
         & GOST R 34.12-2015 Kuznyechik CTR & — & — & Secure & — \\ \cline{2-6}
         & GOST R 34.12-2015 Magma MGM & — & — & — & Insecure \\ \cline{2-6}
         & GOST R 34.12-2015 Kuznyechik MGM & — & — & — & Secure \\ \cline{2-6}
         & IDEA CBC & Insecure & Insecure & Insecure & — \\ \cline{2-6}
         & DES CBC (56) & Insecure & Insecure & Insecure & — \\ \cline{2-6}
         & DES CBC (40) & Insecure & Insecure & — & — \\ \cline{2-6}
         & RC2 CBC & Insecure & Insecure & — & — \\ \hline
         
        \multirow{2}{*}{Stream} & ChaCha20-Poly1305 & — & — & Secure & Secure \\ \cline{2-6}
         & RC4 (128) & Insecure & Insecure & Insecure & — \\ \cline{2-6}
         & RC4 (40) & Insecure & Insecure & — & — \\ \hline
         
        None & Null & Insecure & Insecure & Insecure & — \\ \hline
    \end{tabular}
    \caption{Protocol Version Compatibility of Various Ciphers}
\end{table}


\section{Protocols}

\subsection{TLS 1.2}

TLS 1.2 was released in 2008 and is still considered to be secure, depending on the cipher used. Its biggest change compared to TLS 1.1 is the replacement of the MD5 and SHA-1 hash functions with SHA-256 for all security uses. However TLS 1.2 can still use a MD5 or SHA-1 HMAC for data integrity. TLS 1.2 is still widely used and supported by most modern web browsers and servers. The protocol is defined in RFC 5246.
\\\\
In 2011 refinements defined in RFC 6176 removed the backward compatibility with SSL 2.0.

\subsection{TLS 1.3}

TLS 1.3 was the largest change in the history of TLS, defined in RFC 8446. The TLS 1.3 requirement for forward secrecy meant that a lot of the previously supported key exchange/agreement algorithms were no longer supported. The only key exchange/agreement algorithms supported by TLS 1.3 are different versions of Diffie-Hellman (DH), Elliptic Curve Diffie-Hellman (ECDH) and one algorithm from the Russian GOST standards organization. 
\\\\
With previous TLS versions the encryption and data integrity algorithms were separate. The MAC was computed and appended to the data, then data+MAC were encrypted. However, this has security downsides. TLS 1.3 combined the encryption and data integrity algorithms, only supporting ciphers that use AEAD for data integrity. The supported cipher can be seen in \ref{sec:tls_supported_ciphers}.
\\\\
One of the other major changes, the new handshake protocol, will be discussed in \ref{sec:tls_handshake}.
\\\\
Where as the security of TLS 1.2 depends on the used cipher and key exchange method, TLS 1.3 only supports currently secure methods (expect one of the Russian GOST cithers, which is likely used very little anyway).

\subsubsection{Forward secrecy}

Because forward secrecy, sometimes also called perfect forward secrecy is a major change with TLS 1.3 it is worth explaining what it is. Forward secrecy means that if the long term secrets used for the session, like the private key, are compromised, and the attacker recorded the entire encrypted session, they still can't decrypt the session.
\\\\
Forward secrecy is achieved by using a version of the Diffie Hellman key exchange instead of the client generating a secret key and sending it to the server encrypted with the server's public key. Because the server and client never send the secret key over the network, and the server's private key is no longer a factor in the key generation, an attacker can't decrypt the session if they get the long term secrets.

\section{Subprotocols}

\subsection{Handshake Protocol}

\subsection{Alert Protocol}

\subsection{ChangeCipherSpec Protocol}

\subsection{Application Protocol}



\section{TLS 1.2/1.3 Handshake}
\label{sec:tls_handshake}

\section{Security}

\subsection{Security Features}

\subsection{Attacks}

\section{TLS Record Format}


\begin{table}[htbp]
	\centering
	\caption{TLS record format, general}
	\begin{tabular}{|c|c|c|c|c|}
	\hline
	\textbf{Offset} & \textbf{Byte+0} & \textbf{Byte+1} & \textbf{Byte+2} & \textbf{Byte+3} \\ \hline
	\textbf{Byte 0} & Content type & \multicolumn{3}{c|}{-} \\ \hline
	\multirow{2}{*}{\textbf{Bytes 1-4}} & \multicolumn{2}{c|}{Legacy version} & \multicolumn{2}{c|}{Length} \\ \cline{2-5} 
	 & (Major) & (Minor) & (bits 15-8) & (bits 7-0) \\ \hline
	\textbf{Bytes 5-(\textit{m}-1)} & \multicolumn{4}{c|}{Protocol message(s)} \\ \hline
	\textbf{Bytes \textit{m}-(\textit{p}-1)} & \multicolumn{4}{c|}{MAC (optional)} \\ \hline
	\textbf{Bytes \textit{p}-(\textit{q}-1)} & \multicolumn{4}{c|}{Padding (block ciphers only)} \\ \hline
	\end{tabular}
	\end{table}
	
	
	

\end{document}
